\documentclass[11pt]{article}
\usepackage[utf8]{inputenc}
\usepackage[T1]{fontenc}
\usepackage{amsmath, amssymb}
\usepackage{geometry}
\geometry{a4paper, margin=1in}

\title{Derivation: Why the Simple GIV Isolates Idiosyncratic Shocks}
\author{Gemini CLI}
\date{\today}

\begin{document}

\maketitle

In the \textbf{Simple Model} of Granular Instrumental Variables (Gabaix and Koijen, 2022), the instrument $z_t$ is defined as the difference between size-weighted and equal-weighted outcomes. This document derives why this difference isolates purely idiosyncratic shocks.

\section{The Structural Equations}
The firm-level demand equation is given by:
\begin{equation}
    L_{it} = \phi^d w_t + \eta_t + u_{it}
\end{equation}
where:
\begin{itemize}
    \item $\phi^d w_t$ is the common effect of aggregate wages.
    \item $\eta_t$ is a common macro demand shock hitting all firms equally.
    \item $u_{it}$ is the idiosyncratic shock to firm $i$.
\end{itemize}

\section{Averages and Weight Properties}
The GIV is constructed using two types of averages. A critical property is that both sets of weights sum to 1.

\subsection{Size-Weighted Average ($L_{St}$)}
Defined with weights $S_i$ such that $\sum_{i=1}^N S_i = 1$:
\begin{align}
    L_{St} &= \sum_{i=1}^N S_i (\phi^d w_t + \eta_t + u_{it}) \\
    &= (\phi^d w_t + \eta_t) \left( \sum_{i=1}^N S_i \right) + \sum_{i=1}^N S_i u_{it}
\end{align}
Substituting $\sum S_i = 1$:
\begin{equation}
    L_{St} = (\phi^d w_t + \eta_t) + u_{St}
\end{equation}
where $u_{St} = \sum S_i u_{it}$.

\subsection{Equal-Weighted Average ($L_{Et}$)}
Defined with weights $1/N$ such that $\sum_{i=1}^N \frac{1}{N} = 1$:
\begin{align}
    L_{Et} &= \frac{1}{N} \sum_{i=1}^N (\phi^d w_t + \eta_t + u_{it}) \\
    &= (\phi^d w_t + \eta_t) \left( \frac{1}{N} \sum_{i=1}^N 1 \right) + \frac{1}{N} \sum_{i=1}^N u_{it}
\end{align}
Substituting $\frac{1}{N} \cdot N = 1$:
\begin{equation}
    L_{Et} = (\phi^d w_t + \eta_t) + u_{Et}
\end{equation}
where $u_{Et} = \frac{1}{N} \sum u_{it}$.

\section{The GIV Construction}
The GIV ($z_t$) is defined as the difference:
\begin{align}
    z_t &= L_{St} - L_{Et} \\
    &= \left[ (\phi^d w_t + \eta_t) + u_{St} \right] - \left[ (\phi^d w_t + \eta_t) + u_{Et} \right]
\end{align}
The common components $(\phi^d w_t + \eta_t)$ cancel out perfectly:
\begin{equation}
    z_t = u_{St} - u_{Et}
\end{equation}

\section{Conclusion}
Because both $S_i$ and $1/N$ sum to unity, and because the simple model assumes every firm $i$ has the same loading (implicitly 1) on the macro shock $\eta_t$ and the endogenous wage $w_t$, these aggregate components subtract to zero. The resulting instrument $z_t$ is composed \textit{only} of idiosyncratic shocks, making it a valid instrument for identifying structural elasticities.

\end{document}
